% 康奈尔听课笔记模板（配色取自王道课件）
\documentclass[10.5pt,a4paper,oneside]{article}
\usepackage[a4paper,top=2.3cm,bottom=2.2cm,left=5.7cm,right=1.5cm,
  marginparwidth=4.25cm,marginparsep=0.55cm,headheight=14pt,footskip=1.1cm]{geometry}
\usepackage{xeCJK}
\usepackage{xeCJKfntef}
\xeCJKDeclareCharClass{FullLeft}{"201C,"2018}
\xeCJKDeclareCharClass{FullRight}{"201D,"2019}
\xeCJKDeclareCharClass{CJK}{"2460->"24FF,"2605,"2606}
\setmainfont{Noto Sans CJK SC}[BoldFont={Noto Sans CJK SC Bold},ItalicFont={Noto Sans CJK SC},ItalicFeatures={FakeSlant=0.18}]
\setsansfont{Noto Sans CJK SC}[BoldFont={Noto Sans CJK SC Bold}]
\setmonofont{DejaVu Sans Mono}[Scale=0.9]
\setCJKmainfont{Noto Sans CJK SC}[BoldFont={Noto Sans CJK SC Bold},ItalicFont={Noto Sans CJK SC},ItalicFeatures={FakeSlant=0.18}]
\setCJKsansfont{Noto Sans CJK SC}[BoldFont={Noto Sans CJK SC Bold}]
\setCJKmonofont{Noto Sans Mono CJK SC}[Scale=0.9]
\linespread{1.32}
\usepackage{amsmath,amssymb}
\usepackage[dvipsnames,table]{xcolor}
\usepackage{graphicx}
\usepackage{adjustbox}
\usepackage{tikz}
\usetikzlibrary{calc}
\usepackage[most]{tcolorbox}
\usepackage{eso-pic}
\usepackage{fancyhdr}
\usepackage{titlesec}
\usepackage{needspace}
\usepackage{enumitem}
\usepackage{xltabular,booktabs,array,colortbl}
\usepackage{marginnote}
\reversemarginpar
\renewcommand*{\raggedleftmarginnote}{}
\renewcommand*{\raggedrightmarginnote}{}
\usepackage{caption}
\usepackage{titletoc}
\usepackage{fancyvrb}
\usepackage{framed}
\usepackage[section]{placeins}
\usepackage{hyperref}
\hypersetup{colorlinks=true,linkcolor=wdblue,urlcolor=wdblue,pdftitle={1.2.2
算法的时间复杂度},pdfauthor={听课笔记}}

% ---------- 课件配色 ----------
\definecolor{wdblue}{HTML}{4472C4}
\definecolor{wdlightblue}{HTML}{5B9BD5}
\definecolor{wdtitle}{HTML}{4A86C8}
\definecolor{wddark}{HTML}{1F3864}
\definecolor{wdorange}{HTML}{ED7D31}
\definecolor{wdgray}{HTML}{A5A5A5}
\definecolor{wdgold}{HTML}{FFC000}
\definecolor{wdgreen}{HTML}{70AD47}
\definecolor{wddeepgreen}{HTML}{3F7F2A}
\definecolor{wdred}{HTML}{C00000}
\definecolor{wdpurple}{HTML}{7030A0}
\definecolor{cuebg}{HTML}{F2F6FC}
\definecolor{codebg}{HTML}{F7F9FC}
\color[HTML]{262626}

% ---------- 行内样式 ----------
\newcommand{\hlred}[1]{\textcolor{wdred}{\bfseries #1}}
\newcommand{\hlblue}[1]{\textcolor{wdblue}{\bfseries #1}}
\newcommand{\hlgreen}[1]{\textcolor{wddeepgreen}{\bfseries #1}}
\newcommand{\hlorange}[1]{\textcolor{wdorange}{\bfseries #1}}
\newcommand{\hlpurple}[1]{\textcolor{wdpurple}{\bfseries #1}}
\newcommand{\uln}[1]{\CJKunderline[format=\color{wdred},thickness=0.9pt,depth=0.28em]{#1}}
\newcommand{\wavy}[1]{\CJKunderwave[format=\color{wdorange}]{#1}}
\newcommand{\hlmark}[1]{\CJKunderline[format=\color{wdgold!55},thickness=0.55em,depth=-0.05em]{#1}}
\newcommand{\tagbox}[2]{\tcbox[on line,colback=#1,colframe=#1,coltext=white,boxsep=0.5pt,left=3pt,right=3pt,top=0.5pt,bottom=0.5pt,arc=2pt,fontupper=\small\bfseries]{#2}}

% ---------- 康奈尔版式背景：左侧线索栏 ----------
\newlength{\cuestripw}\setlength{\cuestripw}{5.15cm}
\AddToShipoutPictureBG{%
  \begin{tikzpicture}[remember picture,overlay]
    \fill[cuebg] (current page.north west) rectangle ($(current page.south west)+(\cuestripw,0)$);
    \draw[wdblue,line width=1.2pt] ($(current page.north west)+(\cuestripw,-1.55cm)$) -- ($(current page.south west)+(\cuestripw,1.55cm)$);
    \node[anchor=west,text=wdblue,font=\small\bfseries] at ($(current page.north west)+(0.75cm,-1.25cm)$) {线索 · 问题};
    \node[anchor=west,text=wdblue,font=\small\bfseries] at ($(current page.north west)+(\cuestripw+0.55cm,-1.25cm)$) {笔记 · 老师主线 + 补充讲解};
  \end{tikzpicture}}

% ---------- 页眉页脚 ----------
\pagestyle{fancy}\fancyhf{}
\renewcommand{\headrulewidth}{0pt}
\fancyhead[R]{\small\color{wdgray}第1章 绪论}
\fancyfoot[C]{\small\color{wdgray}\thepage}
\fancyfoot[R]{\small\color{wdorange}1.2.2 算法的时间复杂度 · 听课笔记}
\fancyfoot[L]{\small\hyperlink{toc}{\textcolor{wdblue}{\ensuremath{\uparrow}\,回到目录}}}
\fancypagestyle{plain}{\fancyhf{}\fancyfoot[C]{\small\color{wdgray}\thepage}\fancyfoot[R]{\small\color{wdorange}1.2.2
算法的时间复杂度 · 听课笔记}}

% ---------- 标题：仿课件的蓝色标题 + 六色条 ----------
\newcommand{\sixbar}{\par\nointerlineskip\vspace{3pt}\centering
  \begin{tikzpicture}\foreach \c [count=\i] in {wdlightblue,wdorange,wdgray,wdgold,wdblue,wdgreen}
     {\fill[\c] ({(\i-1)*1.05},0) rectangle ({(\i-1)*1.05+0.92},0.11);}\end{tikzpicture}\par}
\titleformat{\section}[block]{\centering\Large\bfseries\color{wdtitle}}{}{0pt}{}[\sixbar]
\titlespacing*{\section}{0pt}{16pt}{10pt}
\pretocmd{\section}{\Needspace{8\baselineskip}}{}{}
\titleformat{\subsection}{\Needspace{7\baselineskip}\large\bfseries\color{wddark}}{\textcolor{wdblue}{\rule[-0.15em]{4pt}{1.05em}}\hspace{6pt}}{0pt}{}
\titlespacing*{\subsection}{0pt}{12pt}{5pt}
\titleformat{\subsubsection}{\Needspace{5\baselineskip}\normalsize\bfseries\color{wdblue}}{\textcolor{wdorange}{\large\ensuremath{\blacktriangleright}}\hspace{4pt}}{0pt}{}
\titlespacing*{\subsubsection}{0pt}{8pt}{3pt}
\setcounter{secnumdepth}{0}
\setlength{\parindent}{0pt}\setlength{\parskip}{5pt}
\setlist{nosep,leftmargin=1.5em,itemsep=2pt}
\setlist[itemize,1]{label=\textcolor{wdblue}{\small\ensuremath{\bullet}}}
\setlist[itemize,2]{label=\textcolor{wdorange}{\scriptsize\ensuremath{\blacktriangleright}}}
\setlist[enumerate,1]{label=\textcolor{wdblue}{\bfseries\arabic*.}}
\providecommand{\tightlist}{\setlength{\itemsep}{2pt}\setlength{\parskip}{0pt}}

% ---------- 目录 ----------
\setcounter{tocdepth}{2}
\titlecontents{section}[0pt]{\addvspace{5pt}\bfseries\color{wddark}}{}{}{\titlerule*[0.7pc]{\color{wdgray!60}.}\contentspage}
\titlecontents{subsection}[1.6em]{\addvspace{1pt}\small\color{black!75}}{}{}{\titlerule*[0.7pc]{\color{wdgray!40}.}\contentspage}
\newcommand{\tocback}{\par\vspace{-3pt}{\hfill\footnotesize\hyperlink{toc}{\textcolor{wdgray}{\ensuremath{\uparrow}\,回到目录}}}\par\vspace{-4pt}}

% ---------- 线索栏 ----------
\newcommand{\cueblock}[2]{\Needspace{4\baselineskip}\leavevmode#1{#2}\par\vspace{-\baselineskip}\vspace{-\parskip}}
% 线索直接贴在段首行的左侧（不用 marginnote，避免分页时偶尔落进正文栏）
\newcommand{\cuebox}[1]{\leavevmode\llap{\smash{\begin{minipage}[t]{\marginparwidth}\raggedright\small\linespread{1.18}\selectfont #1\end{minipage}}\hspace{\marginparsep}}}
\newcommand{\cuenote}[1]{\cuebox{{\color{wdblue}\bfseries Q\,}\color{wddark}#1}}
\newcommand{\cuekey}[1]{\cuebox{\color{wdred}\bfseries #1}}

% ---------- 讲解框 ----------
\tcbset{wdbox/.style={enhanced,breakable,arc=2pt,boxrule=0pt,leftrule=3pt,left=7pt,right=7pt,top=5pt,bottom=5pt,
   before skip=7pt,after skip=7pt,fonttitle=\bfseries\small,attach title to upper={\quad},coltitle=#1,
   colframe=#1,colback=#1!6}}
\newtcolorbox{defbox}[1][]{enhanced,breakable,arc=3pt,boxrule=0.8pt,colframe=wdblue,colback=wdblue!4,
   left=8pt,right=8pt,top=6pt,bottom=6pt,fonttitle=\bfseries\small,coltitle=white,colbacktitle=wdblue,
   title={定义#1},before skip=8pt,after skip=8pt}
\newtcolorbox{bridgebox}{wdbox=wdblue,title={承上启下}}
\newtcolorbox{plainbox}{wdbox=wddeepgreen,colframe=wdgreen,colback=wdgreen!8,title={大白话}}
\newtcolorbox{deepbox}{wdbox=wdpurple,title={背后的思想}}
\newtcolorbox{exambox}{wdbox=wdred,title={考点}}
\newtcolorbox{pitfallbox}{wdbox=wdorange,title={易错提醒}}
\newtcolorbox{bookbox}{wdbox=wddark,colback=wdgray!10,colframe=wdgray,title={王道书补充}}
\newtcolorbox{teacherbox}{wdbox=wdorange,colback=wdgold!8,colframe=wdgold,coltitle=wdorange,title={老师原话}}
\newtcolorbox{fixbox}{wdbox=wdred,colback=wdred!4,title={勘误 / 说明}}
% 习题笔记用
\newtcolorbox{questionbox}[1][]{enhanced,breakable,arc=3pt,boxrule=0.8pt,colframe=wddark,colback=white,
   left=8pt,right=8pt,top=6pt,bottom=6pt,fonttitle=\bfseries\small,coltitle=white,colbacktitle=wddark,
   title={#1},before skip=10pt,after skip=6pt}
\newtcolorbox{videobox}{wdbox=wdorange,colback=wdgold!8,colframe=wdgold,coltitle=wdorange,title={视频讲解}}
\newtcolorbox{jiexibox}{wdbox=wddark,colback=wdgray!10,colframe=wdgray,title={书上解析}}
\newtcolorbox{addbox}{wdbox=wdpurple,title={我的补充}}
\newtcolorbox{recallbox}{wdbox=wdblue,title={回顾正课}}
\newtcolorbox{watchbox}{enhanced,breakable,arc=3pt,boxrule=1pt,colframe=wdorange,colback=wdgold!12,
   fonttitle=\bfseries\small,coltitle=white,colbacktitle=wdorange,title={建议去看视频},left=8pt,right=8pt,before skip=8pt,after skip=8pt}
\newtcolorbox{summarybox}{enhanced,arc=3pt,
   boxrule=1pt,colframe=wdblue,colback=white,colbacktitle=wdblue,coltitle=white,fonttitle=\bfseries,
   title={总结区 · 用自己的话复述这一节},left=10pt,right=10pt,top=8pt,bottom=8pt,before skip=14pt}

% ---------- 代码 ----------
\usepackage{color}
\usepackage{fancyvrb}
\newcommand{\VerbBar}{|}
\newcommand{\VERB}{\Verb[commandchars=\\\{\}]}
\DefineVerbatimEnvironment{Highlighting}{Verbatim}{commandchars=\\\{\}}
% Add ',fontsize=\small' for more characters per line
\newenvironment{Shaded}{}{}
\newcommand{\AlertTok}[1]{\textcolor[rgb]{1.00,0.00,0.00}{#1}}
\newcommand{\AnnotationTok}[1]{\textcolor[rgb]{0.38,0.63,0.69}{#1}}
\newcommand{\AttributeTok}[1]{\textcolor[rgb]{0.49,0.56,0.16}{#1}}
\newcommand{\BaseNTok}[1]{\textcolor[rgb]{0.77,0.35,0.07}{#1}}
\newcommand{\BuiltInTok}[1]{\textcolor[rgb]{0.00,0.50,0.00}{#1}}
\newcommand{\CharTok}[1]{\textcolor[rgb]{0.63,0.19,0.43}{#1}}
\newcommand{\CommentTok}[1]{\textcolor[rgb]{0.18,0.55,0.23}{#1}}
\newcommand{\CommentVarTok}[1]{\textcolor[rgb]{0.18,0.55,0.23}{#1}}
\newcommand{\ConstantTok}[1]{\textcolor[rgb]{0.77,0.35,0.07}{#1}}
\newcommand{\ControlFlowTok}[1]{\textcolor[rgb]{0.12,0.31,0.71}{#1}}
\newcommand{\DataTypeTok}[1]{\textcolor[rgb]{0.12,0.31,0.71}{#1}}
\newcommand{\DecValTok}[1]{\textcolor[rgb]{0.77,0.35,0.07}{#1}}
\newcommand{\DocumentationTok}[1]{\textcolor[rgb]{0.18,0.55,0.23}{#1}}
\newcommand{\ErrorTok}[1]{\textcolor[rgb]{1.00,0.00,0.00}{#1}}
\newcommand{\ExtensionTok}[1]{#1}
\newcommand{\FloatTok}[1]{\textcolor[rgb]{0.77,0.35,0.07}{#1}}
\newcommand{\FunctionTok}[1]{\textcolor[rgb]{0.44,0.19,0.63}{#1}}
\newcommand{\ImportTok}[1]{\textcolor[rgb]{0.44,0.19,0.63}{#1}}
\newcommand{\InformationTok}[1]{\textcolor[rgb]{0.38,0.63,0.69}{#1}}
\newcommand{\KeywordTok}[1]{\textcolor[rgb]{0.12,0.31,0.71}{#1}}
\newcommand{\NormalTok}[1]{#1}
\newcommand{\OperatorTok}[1]{\textcolor[rgb]{0.20,0.20,0.20}{#1}}
\newcommand{\OtherTok}[1]{\textcolor[rgb]{0.00,0.44,0.13}{#1}}
\newcommand{\PreprocessorTok}[1]{\textcolor[rgb]{0.44,0.19,0.63}{#1}}
\newcommand{\RegionMarkerTok}[1]{#1}
\newcommand{\SpecialCharTok}[1]{\textcolor[rgb]{0.63,0.19,0.43}{#1}}
\newcommand{\SpecialStringTok}[1]{\textcolor[rgb]{0.63,0.19,0.43}{#1}}
\newcommand{\StringTok}[1]{\textcolor[rgb]{0.63,0.19,0.43}{#1}}
\newcommand{\VariableTok}[1]{\textcolor[rgb]{0.10,0.09,0.49}{#1}}
\newcommand{\VerbatimStringTok}[1]{\textcolor[rgb]{0.63,0.19,0.43}{#1}}
\newcommand{\WarningTok}[1]{\textcolor[rgb]{0.38,0.63,0.69}{#1}}
\makeatletter
\@ifundefined{Shaded}{}{\renewenvironment{Shaded}{\def\FrameCommand{\fboxsep=5pt\fboxrule=0.5pt\fcolorbox{wdblue!35}{codebg}}%
   \MakeFramed{\advance\hsize-\width\FrameRestore}\small}{\endMakeFramed}}
\makeatother
\fvset{fontsize=\small}

% ---------- 表格 ----------
\renewcommand{\arraystretch}{1.42}
\newcolumntype{Y}{>{\raggedright\arraybackslash}X}

% ---------- 图片 ----------
% 课件截图做成浮动体：LaTeX 会自动用正文填满页面，避免大块空白
\setcounter{topnumber}{3}\setcounter{bottomnumber}{2}\setcounter{totalnumber}{4}
\renewcommand{\topfraction}{0.85}\renewcommand{\bottomfraction}{0.6}\renewcommand{\textfraction}{0.12}\renewcommand{\floatpagefraction}{0.75}
\setlength{\textfloatsep}{10pt plus 3pt minus 2pt}\setlength{\floatsep}{8pt plus 2pt}
\newcommand{\slidefig}[3]{\begin{figure}[!ht]\centering\begin{tcolorbox}[enhanced,width=#2\linewidth,arc=2pt,boxrule=0.6pt,colframe=wdgray!70,
   colback=white,left=2pt,right=2pt,top=2pt,bottom=2pt,drop fuzzy shadow=black!25]\includegraphics[width=\linewidth]{#1}\end{tcolorbox}
   \vspace{-3pt}{\small\color{wdgray}#3}\end{figure}}
\newcommand{\plainfig}[3]{\par\noindent\begin{minipage}{\linewidth}\centering\includegraphics[width=#2\linewidth]{#1}\par\vspace{2pt}{\small\color{wdgray}#3}\end{minipage}\par\vspace{4pt}}

\begin{document}
% ---------- 封面条（仿课件封面） ----------
\thispagestyle{plain}
\begin{tcolorbox}[enhanced,enlarge left by=-4.8cm,width=\linewidth+4.8cm,arc=0pt,boxrule=0pt,
   colback=wddark,interior style={left color=wddark,right color=wdblue},left=18pt,right=18pt,top=14pt,bottom=12pt,
   overlay={\draw[white,line width=2.5pt] ([xshift=10pt,yshift=-10pt]frame.north west) -- ++(0,-0.9cm);
            \draw[white,line width=2.5pt] ([xshift=10pt,yshift=-10pt]frame.north west) -- ++(1.6cm,0);
            \draw[white,line width=2.5pt] ([xshift=-10pt,yshift=10pt]frame.south east) -- ++(0,0.9cm);
            \draw[white,line width=2.5pt] ([xshift=-10pt,yshift=10pt]frame.south east) -- ++(-1.6cm,0);}]
  {\color{white!85!wdlightblue}\small 数据结构 · 第1章
绪论 · 康奈尔听课笔记}\par\vspace{6pt}
  {\color{white}\fontsize{22}{28}\selectfont\bfseries 1.2.2
算法的时间复杂度}\par\vspace{6pt}
  {\color{white!80!wdlightblue}\small 课件 1.2\_2 · 字幕 P6 · 王道书
§1.2.2（时间复杂度）+ 本章“归纳总结”}
\end{tcolorbox}
\vspace{4pt}
\begin{tcolorbox}[enhanced,arc=3pt,boxrule=0.8pt,colframe=wdblue,colback=wdblue!3,left=10pt,right=10pt,top=6pt,bottom=6pt,
   colbacktitle=wdblue,coltitle=white,fonttitle=\bfseries\small,title={目录 · 点击标题跳转}]
\hypertarget{toc}{}\makeatletter\@starttoc{toc}\makeatother
\end{tcolorbox}
\hypertarget{ux672cux8282ux5b9aux4f4dux7b2c-1-ux7ae0ux771fux6b63ux7684ux8003ux70b9}{%
\section{本节定位：第 1
章真正的考点}\label{ux672cux8282ux5b9aux4f4dux7b2c-1-ux7ae0ux771fux6b63ux7684ux8003ux70b9}}

\tocback

\cueblock{\cuenote}{\relax{时间复杂度在考试里有多重要？}}

\begin{exambox}

王道书考点追踪：\textbf{分析算法的时间复杂度}几乎年年考（2011---2014、2017、2019、2022、2025
年选择题；算法设计大题也几乎每年要求写时间复杂度）。这是整个第 1
章分量最重的一节。

\end{exambox}

\begin{figure}[!ht]\centering\adjustimage{scale=0.85,max width=\linewidth,max height=0.6\textheight}{/mnt/project-files/notes/1数据结构/笔记/_build/cache/64e789639cc34691ec27fe71104ded3417d23b63.pdf}\end{figure}

\hypertarget{ux4e00ux600eux4e48ux8bc4ux4f30ux4e00ux4e2aux7b97ux6cd5ux7684ux65f6ux95f4ux5f00ux9500}{%
\section{一、怎么评估一个算法的时间开销？}\label{ux4e00ux600eux4e48ux8bc4ux4f30ux4e00ux4e2aux7b97ux6cd5ux7684ux65f6ux95f4ux5f00ux9500}}

\tocback

\hypertarget{ux5148ux8dd1ux4e00ux904dux518dux8ba1ux65f6ux4e8bux540eux7edfux8ba1}{%
\subsection{1.
先跑一遍再计时？（事后统计）}\label{ux5148ux8dd1ux4e00ux904dux518dux8ba1ux65f6ux4e8bux540eux7edfux8ba1}}

\Needspace{4\baselineskip}

\cuenote{“事后统计”有哪四个问题？}最直觉的想法：把算法写成代码，让机器跑一下，看运行了多久。老师指出这样做有四个问题：

\slidefig{assets/1.2.2-p03.jpg}{0.65}{课件：如何评估算法时间开销？}

{\rowcolors{2}{wdblue!5}{white}\arrayrulecolor{wdblue!35}
\begin{xltabular}{\linewidth}{>{\hsize=1.048\hsize\raggedright\arraybackslash}X>{\hsize=0.952\hsize\raggedright\arraybackslash}X}
\rowcolor{wdblue}\textcolor{white}{\bfseries 问题} & \textcolor{white}{\bfseries 说明} \\
\endfirsthead
① 和\textbf{机器性能}有关 & 超级计算机 vs 单片机，同一个算法跑出来的时间天差地别 \\ \hline
② 和\textbf{编程语言}有关 & 越高级的语言执行效率越低，同一算法用 C 写一般比用 Java 写快 \\ \hline
③ 和\textbf{编译程序产生的机器指令质量}有关 & 编译器优化得好不好也会影响时间（看不懂没关系） \\ \hline
④ 有些算法\textbf{不能事后再统计} & 比如导弹控制算法：等导弹打错地方再统计时间，已经没有意义了 \\ \hline
\end{xltabular}\par\vspace{4pt}}

①②③ 说明要排除与算法无关的外界因素，④
说明要事先估计，于是有了时间复杂度：\hlred{事前预估时间开销 \(T(n)\) 与问题规模 \(n\) 的关系}。

\hypertarget{ux95eeux9898ux89c4ux6a21-n-ux662fux4ec0ux4e48}{%
\subsection{\texorpdfstring{2. 问题规模 \(n\)
是什么}{2. 问题规模 n 是什么}}\label{ux95eeux9898ux89c4ux6a21-n-ux662fux4ec0ux4e48}}

\Needspace{4\baselineskip}

\cuenote{为什么时间开销要写成 \(n\) 的函数？}问题规模越大，运行时间理应越长，所以时间开销
\(T\) 应表示成问题规模 \(n\)
的函数，分析时间复杂度的关键就是\uln{用 \(n\) 表示出 \(T\)}。

\hypertarget{ux4e8cux7528ux7b97ux6cd5ux8868ux767dux7231ux4f60-n-ux904d}{%
\section{\texorpdfstring{二、用算法表白：“爱你 n
遍”}{二、用算法表白：爱你 n 遍}}\label{ux4e8cux7528ux7b97ux6cd5ux8868ux767dux7231ux4f60-n-ux904d}}

\tocback

\Needspace{4\baselineskip}

\cuenote{算法 1 中每一行分别执行了几次？为什么 while 的判断是 3001 次？}老师从电影《复仇者联盟
4》里钢铁侠和女儿的对话（“爱你三千遍”）出发，写了一个“逐步递增型爱你”：

\slidefig{assets/1.2.2-p05.jpg}{0.8}{课件：算法 1 与语句频度}

\begin{Shaded}
\begin{Highlighting}[]
\CommentTok{//算法1— 逐步递增型爱你}
\DataTypeTok{void}\NormalTok{ loveYou}\OperatorTok{(}\DataTypeTok{int}\NormalTok{ n}\OperatorTok{)} \OperatorTok{\{}         \CommentTok{//n 为问题规模}
    \DataTypeTok{int}\NormalTok{ i}\OperatorTok{=}\DecValTok{1}\OperatorTok{;}                  \CommentTok{// ① 爱你的程度}
    \ControlFlowTok{while}\OperatorTok{(}\NormalTok{i}\OperatorTok{\textless{}=}\NormalTok{n}\OperatorTok{)\{}              \CommentTok{// ②}
\NormalTok{        i}\OperatorTok{++;}                  \CommentTok{// ③ 每次+1}
\NormalTok{        printf}\OperatorTok{(}\StringTok{"I Love You }\SpecialCharTok{\%d\textbackslash{}n}\StringTok{"}\OperatorTok{,}\NormalTok{ i}\OperatorTok{);}    \CommentTok{// ④}
    \OperatorTok{\}}
\NormalTok{    printf}\OperatorTok{(}\StringTok{"I Love You More Than }\SpecialCharTok{\%d\textbackslash{}n}\StringTok{"}\OperatorTok{,}\NormalTok{ n}\OperatorTok{);}  \CommentTok{// ⑤}
\OperatorTok{\}}

\DataTypeTok{int}\NormalTok{ main}\OperatorTok{()\{}
\NormalTok{    loveYou}\OperatorTok{(}\DecValTok{3000}\OperatorTok{);}
\OperatorTok{\}}
\end{Highlighting}
\end{Shaded}

\textbf{先做一个简化}：假设执行任意一行代码所花的时间都相同（实际不同，但考虑太多因素就没法评价了）。那么只要数出每行执行了几次就行。

\begin{fixbox}

\textbf{名词统一}：课件写“语句频度”，老师口头说“执行次数”，王道书写“频度”，三者是\uln{同一个意思：某条语句被执行的次数}。本笔记下文统一说“执行次数”。

\end{fixbox}

{\rowcolors{2}{wdblue!5}{white}\arrayrulecolor{wdblue!35}
\begin{xltabular}{\linewidth}{>{\hsize=0.360\hsize\raggedright\arraybackslash}X>{\hsize=0.840\hsize\raggedright\arraybackslash}X>{\hsize=1.800\hsize\raggedright\arraybackslash}X}
\rowcolor{wdblue}\textcolor{white}{\bfseries 语句} & \textcolor{white}{\bfseries 执行次数（n = 3000）} & \textcolor{white}{\bfseries 为什么} \\
\endfirsthead
① & 1 次 & 只在开头执行一次 \\ \hline
② & \hlred{3001 次} & 循环体执行 3000 次，每次执行前都要判断一次；第 3000 次循环后 \texttt{i} 变成 3001，还要\textbf{再判断一次}发现 \texttt{3001\textless{}=3000} 不成立，才跳出循环。所以判断比循环体多 1 次 \\ \hline
③④ & 各 3000 次 & 循环体执行了 3000 次 \\ \hline
⑤ & 1 次 & 循环结束后执行一次 \\ \hline
\end{xltabular}\par\vspace{4pt}}

\[T(3000) = 1 + 3001 + 2\times 3000 + 1\]

把 3000 换成 \(n\)，得到时间开销与问题规模的关系：

\[T(n) = 1 + (n+1) + 2n + 1 = 3n + 3\]

\begin{pitfallbox}

“\texttt{while}
的条件判断比循环体多执行一次”这个细节，问“某条语句执行了几次”时会用到。\texttt{for(i=1;\ i\textless{}=n;\ i++)}
也一样：\texttt{i\textless{}=n} 判断 \(n+1\) 次，循环体执行 \(n\) 次。

\end{pitfallbox}

由此引出两个问题：\textbf{问题 1} 表达式能否简化？\textbf{问题 2}
几千行代码也要逐行数吗？

\hypertarget{ux4e09ux95eeux9898-1ux8868ux8fbeux5f0fux80fdux4e0dux80fdux7b80ux5316ux5927-o-ux8868ux793aux6cd5}{%
\section{三、问题 1：表达式能不能简化？------大 O
表示法}\label{ux4e09ux95eeux9898-1ux8868ux8fbeux5f0fux80fdux4e0dux80fdux7b80ux5316ux5927-o-ux8868ux793aux6cd5}}

\tocback

\hypertarget{ux53eaux4fddux7559ux6700ux9ad8ux9636ux518dux53bbux6389ux7cfbux6570}{%
\subsection{1.
只保留最高阶，再去掉系数}\label{ux53eaux4fddux7559ux6700ux9ad8ux9636ux518dux53bbux6389ux7cfbux6570}}

\Needspace{4\baselineskip}

\cuenote{为什么可以只看最高阶项？为什么系数也能忽略？}假设有三个算法：

\[T_1(n)=3n+3,\qquad T_2(n)=n^2+3n+1000,\qquad T_3(n)=n^3+n^2+9999999\]

\slidefig{assets/1.2.2-p06.jpg}{0.8}{课件：当问题规模 n 足够大时}

取 \(n=3000\)，比较“只保留最高阶项”和“精确值”：

{\rowcolors{2}{wdblue!5}{white}\arrayrulecolor{wdblue!35}
\begin{xltabular}{\linewidth}{>{\hsize=0.714\hsize\raggedright\arraybackslash}X>{\hsize=1.143\hsize\raggedright\arraybackslash}X>{\hsize=1.143\hsize\raggedright\arraybackslash}X}
\rowcolor{wdblue}\textcolor{white}{\bfseries } & \textcolor{white}{\bfseries 只保留最高阶项} & \textcolor{white}{\bfseries 精确值} \\
\endfirsthead
\(T_1\) & \(3n = 9000\) & \(9003\) \\ \hline
\(T_2\) & \(n^2 = 9{,}000{,}000\) & \(9{,}010{,}000\) \\ \hline
\(T_3\) & \(n^3 = 27{,}000{,}000{,}000\) & \(27{,}018{,}999{,}999\) \\ \hline
\end{xltabular}\par\vspace{4pt}}

差别都很小。\textbf{结论 1}：\hlred{可以只考虑阶数高的部分}。

\begin{teacherbox}

你有 3000 块，舍友有 9000 块，谁也不比谁富多少，是同一个数量级；有人有
900 万，那才不是一个数量级。

\end{teacherbox}

再看一组对比：

\begin{itemize}
\tightlist
\item
  当 \(n=3000\) 时，\(9999n = 29{,}997{,}000\)，\uln{远小于}
  \(n^3 = 27{,}018{,}999{,}999\)；
\item
  当 \(n=1000000\) 时，\(9999n = 9{,}999{,}000{,}000\)，\uln{远小于}
  \(n^2 = 1{,}000{,}000{,}000{,}000\)。
\end{itemize}

\textbf{结论
2}：\hlred{问题规模足够大时，常数项系数也可以忽略}。哪怕系数是
9999，只要 \(n\) 够大，低阶终究打不过高阶。

\hypertarget{ux5927-o-ux8bb0ux6cd5}{%
\subsection{2. 大 O 记法}\label{ux5927-o-ux8bb0ux6cd5}}

\cueblock{\cuenote}{\relax{大 O 的数学含义？}}

\begin{defbox}[：大 O 表示“同阶”]

\(T(n)=O(f(n))\) 表示 \(T(n)\) 与 \(f(n)\) 同阶、同等数量级。即当
\(n\to\infty\) 时，二者之比为常数：

\[\lim_{n\to\infty}\frac{T(n)}{f(n)} = k\quad(k\ \text{为常数})\]

\end{defbox}

于是：\(T_1(n)=O(n)\)，\(T_2(n)=O(n^2)\)，\(T_3(n)=O(n^3)\)。

\textbf{操作口诀}：\uln{只保留最高阶项 → 系数变为 1 → 外面套上 \(O(\ )\)}。

\begin{bookbox}

王道书给出的是更严格的定义：若存在正常数 \(C\) 和 \(n_0\)，使得对所有
\(n\ge n_0\)，都有 \(0\le T(n)\le C f(n)\)，则记
\(T(n)=O(f(n))\)。它表示 \textbf{当 \(n\) 足够大时，\(T(n)\)
的增长速度不超过 \(f(n)\) 的常数倍}。

书上还说：时间复杂度通常由算法中\textbf{基本运算（最深层循环内的关键操作）}的执行次数决定，该执行次数的数量级即为整个算法的时间复杂度。

\end{bookbox}

\begin{deepbox}

课上说“同阶”，书上的定义只要求“增长不超过 \(f(n)\) 的常数倍”，所以严格说
\(3n+3=O(n^2)\) 也成立。但做题一律写\textbf{最紧的上界}：\(3n+3\) 写
\(O(n)\)。

\end{deepbox}

\hypertarget{ux52a0ux6cd5ux89c4ux5219ux4e0eux4e58ux6cd5ux89c4ux5219}{%
\subsection{3.
加法规则与乘法规则}\label{ux52a0ux6cd5ux89c4ux5219ux4e0eux4e58ux6cd5ux89c4ux5219}}

\cueblock{\cuenote}{\relax{加法规则、乘法规则分别是什么？各在什么代码结构里用？}}

\slidefig{assets/1.2.2-p07.jpg}{0.8}{课件：加法规则与乘法规则}

\begin{defbox}

\textbf{a）加法规则}：\(T(n)=T_1(n)+T_2(n)=O(f(n))+O(g(n))=O(\max(f(n),g(n)))\)
\hlred{多项相加，只保留最高阶的项，且系数变为 1}。

\textbf{b）乘法规则}：\(T(n)=T_1(n)\times T_2(n)=O(f(n))\times O(g(n))=O(f(n)\times g(n))\)
\hlred{多项相乘，都保留}。

\end{defbox}

\begin{bookbox}

王道书把两个规则和代码结构对应起来，这样就知道“什么时候用哪个”：

\begin{itemize}
\tightlist
\item
  \textbf{加法规则}：两段代码\textbf{顺序执行}，总复杂度取两者中的高阶项。
\item
  \textbf{乘法规则}：一段代码\textbf{嵌套在另一段内部}，总复杂度为两者之积。
\end{itemize}

例：设语句块 \texttt{a\{\}}、\texttt{b\{\}}、\texttt{c\{\}}
的时间复杂度分别为 \(O(1)\)、\(O(n)\)、\(O(n^2)\)：

\begin{Shaded}
\begin{Highlighting}[]
\NormalTok{a}\OperatorTok{\{}\NormalTok{ b}\OperatorTok{\{\}}\NormalTok{ c}\OperatorTok{\{\}} \OperatorTok{\}}     \CommentTok{// b、c 并列在 a 里：O(n) + O(n\^{}2) = O(n\^{}2)，加法规则}
\NormalTok{a}\OperatorTok{\{}\NormalTok{ b}\OperatorTok{\{}\NormalTok{ c}\OperatorTok{\{\}} \OperatorTok{\}} \OperatorTok{\}}    \CommentTok{// c 嵌套在 b 里：O(n) × O(n\^{}2) = O(n\^{}3)，乘法规则}
\end{Highlighting}
\end{Shaded}

\end{bookbox}

\hypertarget{ux8c01ux7684ux9636ux66f4ux9ad8ux5e38ux5bf9ux5e42ux6307ux9636}{%
\subsection{\texorpdfstring{4.
谁的阶更高？------“常对幂指阶”}{4. 谁的阶更高？------常对幂指阶}}\label{ux8c01ux7684ux9636ux66f4ux9ad8ux5e38ux5bf9ux5e42ux6307ux9636}}

\Needspace{4\baselineskip}

\cuenote{常见数量级从小到大怎么排？口诀？}用加法规则时必须判断“谁的阶更高”。课件例子：

\[T_3(n)=n^3+n^2\log_2 n = O(n^3)+O(n^2\log_2 n) = \ ???\]

把 \(n^3\) 写成 \(n^2\cdot n\)，两项都有 \(n^2\)，于是只需比较 \(n\) 和
\(\log_2 n\) 谁增长得快：

\[\lim_{n\to\infty}\frac{\log_2 n}{n}\ \overset{\text{洛必达}}{=}\ \lim_{n\to\infty}\frac{1}{n\ln 2}=0\]

极限为 0，说明 \(n\to\infty\) 时 \textbf{\(n\) 比 \(\log_2 n\)
变大的速度快很多}，所以 \(T_3(n)=O(n^3)\)。

\slidefig{assets/1.2.2-p08.jpg}{0.8}{课件：常对幂指阶}

同理可证 \(2^n\) 比 \(n^2\) 增长得快：

\[\lim_{n\to\infty}\frac{n^2}{2^n}\ \overset{\text{洛必达}}{=}\ \lim_{n\to\infty}\frac{2n}{2^n\ln 2}\ \overset{\text{洛必达}}{=}\ \lim_{n\to\infty}\frac{2}{2^n\ln^2 2}=0\]

\begin{defbox}[：常见数量级（必须背熟）]

\[O(1)<O(\log_2 n)<O(n)<O(n\log_2 n)<O(n^2)<O(n^3)<O(2^n)<O(n!)<O(n^n)\]

口诀：\hlred{\textbf{常对幂指阶}}（常数 \textless{} 对数 \textless{}
幂函数 \textless{} 指数 \textless{}
阶乘），老师称之为“巴拉拉能量------常对幂指阶”。

\end{defbox}

\slidefig{assets/1.2.2-p09.jpg}{0.62}{课件：各数量级的增长曲线}

\begin{pitfallbox}

\begin{itemize}
\tightlist
\item
  不要求会用极限证明，\textbf{记住顺序}就够用。
\item
  \(n\log_2 n\) 夹在 \(n\) 和 \(n^2\)
  之间，是排序算法里最常见的“好”复杂度（第 8 章）。
\item
  对数的底数在大 O 里不重要（\(\log_2 n\) 和 \(\log_3 n\)
  只差常数倍），但课件和真题习惯写 \(\log_2 n\)。
\end{itemize}

\end{pitfallbox}

\hypertarget{ux56dbux95eeux9898-2ux51e0ux5343ux884cux4ee3ux7801ux4e5fux8981ux4e00ux884cux884cux6570ux5417}{%
\section{四、问题
2：几千行代码也要一行行数吗？}\label{ux56dbux95eeux9898-2ux51e0ux5343ux884cux4ee3ux7801ux4e5fux8981ux4e00ux884cux884cux6570ux5417}}

\tocback

\hypertarget{ux7ed3ux8bba-1ux987aux5e8fux6267ux884cux7684ux4ee3ux7801ux53eaux5f71ux54cdux5e38ux6570ux9879}{%
\subsection{结论
1：顺序执行的代码只影响常数项}\label{ux7ed3ux8bba-1ux987aux5e8fux6267ux884cux7684ux4ee3ux7801ux53eaux5f71ux54cdux5e38ux6570ux9879}}

\cueblock{\cuenote}{\relax{顺序执行的代码要不要数？}}

\slidefig{assets/1.2.2-p11.jpg}{0.8}{课件：在开头插入 1000 行顺序执行的代码}

在算法 1 前面插入 1000 行\textbf{顺序执行}（不带循环）的代码：

\[T(3000)=1000+1+3001+2\times3000+1,\qquad T(n)=3n+1003=O(n)\]

常数从 3 变成 1003，但数量级仍是
\(O(n)\)。\hlred{顺序执行的代码只会影响常数项，可以忽略}。

\hypertarget{ux7ed3ux8bba-2ux53eaux9700ux6311ux5faaux73afux4e2dux7684ux4e00ux4e2aux57faux672cux64cdux4f5c}{%
\subsection{结论
2：只需挑循环中的一个基本操作}\label{ux7ed3ux8bba-2ux53eaux9700ux6311ux5faaux73afux4e2dux7684ux4e00ux4e2aux57faux672cux64cdux4f5c}}

\Needspace{4\baselineskip}

\cuenote{循环里的三条语句都要数吗？}之前分析循环时，把 ②③④
都数进去了，得到 \(3n\)。但最后系数总要变成
1，所以\hlred{只需挑循环中的一个基本操作，分析它的执行次数与 \(n\) 的关系即可}。比如只看语句
④，执行 \(n\) 次，\(T(n)=O(n)\)。

\hypertarget{ux7ed3ux8bba-3ux591aux5c42ux5d4cux5957ux5faaux73afux53eaux770bux6700ux6df1ux5c42}{%
\subsection{结论
3：多层嵌套循环只看最深层}\label{ux7ed3ux8bba-3ux591aux5c42ux5d4cux5957ux5faaux73afux53eaux770bux6700ux6df1ux5c42}}

\cueblock{\cuenote}{\relax{嵌套循环怎么分析？}}

\begin{Shaded}
\begin{Highlighting}[]
\CommentTok{//算法2— 嵌套循环型爱你}
\DataTypeTok{void}\NormalTok{ loveYou}\OperatorTok{(}\DataTypeTok{int}\NormalTok{ n}\OperatorTok{)} \OperatorTok{\{}          \CommentTok{//n 为问题规模}
    \DataTypeTok{int}\NormalTok{ i}\OperatorTok{=}\DecValTok{1}\OperatorTok{;}                   \CommentTok{//爱你的程度}
    \ControlFlowTok{while}\OperatorTok{(}\NormalTok{i}\OperatorTok{\textless{}=}\NormalTok{n}\OperatorTok{)\{}               \CommentTok{// 外层循环执行 n 次}
\NormalTok{        i}\OperatorTok{++;}                   \CommentTok{//每次+1}
\NormalTok{        printf}\OperatorTok{(}\StringTok{"I Love You }\SpecialCharTok{\%d\textbackslash{}n}\StringTok{"}\OperatorTok{,}\NormalTok{ i}\OperatorTok{);}
        \ControlFlowTok{for} \OperatorTok{(}\DataTypeTok{int}\NormalTok{ j}\OperatorTok{=}\DecValTok{1}\OperatorTok{;}\NormalTok{ j}\OperatorTok{\textless{}=}\NormalTok{n}\OperatorTok{;}\NormalTok{ j}\OperatorTok{++)\{}        \CommentTok{// 嵌套两层循环}
\NormalTok{            printf}\OperatorTok{(}\StringTok{"I am Iron Man}\SpecialCharTok{\textbackslash{}n}\StringTok{"}\OperatorTok{);}   \CommentTok{// 内层循环共执行 n\^{}2 次}
        \OperatorTok{\}}
    \OperatorTok{\}}
\NormalTok{    printf}\OperatorTok{(}\StringTok{"I Love You More Than }\SpecialCharTok{\%d\textbackslash{}n}\StringTok{"}\OperatorTok{,}\NormalTok{ n}\OperatorTok{);}
\OperatorTok{\}}
\end{Highlighting}
\end{Shaded}

\slidefig{assets/1.2.2-p12.jpg}{0.8}{课件：算法 2 与三个结论}

外层每执行 1 次，内层就执行 \(n\) 次；外层执行 \(n\) 次，内层一共执行
\(n\times n=n^2\) 次：

\[T(n) = O(n) + O(n^2) = O(n^2)\]

\hlred{如果有多层嵌套循环，只需关注最深层循环循环了几次}。

\begin{deepbox}

\textbf{三步法}（课件知识回顾）：① 找最深层循环里的一个基本操作；②
求它的执行次数 \(x=f(n)\)；③ \(O(x)\) 即
\(T(n)\)。最深层执行次数最多，所以加法规则下只留它。

\end{deepbox}

\hypertarget{ux4e94ux4e24ux4e2aux5c0fux7ec3ux4e60}{%
\section{五、两个小练习}\label{ux4e94ux4e24ux4e2aux5c0fux7ec3ux4e60}}

\tocback

\hypertarget{ux5c0fux7ec3ux4e60-1ux6307ux6570ux9012ux589eux578bux7231ux4f60}{%
\subsection{小练习
1：指数递增型爱你}\label{ux5c0fux7ec3ux4e60-1ux6307ux6570ux9012ux589eux578bux7231ux4f60}}

\cueblock{\cuenote}{\relax{循环变量每次翻倍时，循环执行几次？}}

\begin{Shaded}
\begin{Highlighting}[]
\CommentTok{//算法3— 指数递增型爱你}
\DataTypeTok{void}\NormalTok{ loveYou}\OperatorTok{(}\DataTypeTok{int}\NormalTok{ n}\OperatorTok{)} \OperatorTok{\{}      \CommentTok{//n 为问题规模}
    \DataTypeTok{int}\NormalTok{ i}\OperatorTok{=}\DecValTok{1}\OperatorTok{;}               \CommentTok{//爱你的程度}
    \ControlFlowTok{while}\OperatorTok{(}\NormalTok{i}\OperatorTok{\textless{}=}\NormalTok{n}\OperatorTok{)\{}
\NormalTok{        i}\OperatorTok{=}\NormalTok{i}\OperatorTok{*}\DecValTok{2}\OperatorTok{;}             \CommentTok{//每次翻倍}
\NormalTok{        printf}\OperatorTok{(}\StringTok{"I Love You }\SpecialCharTok{\%d\textbackslash{}n}\StringTok{"}\OperatorTok{,}\NormalTok{ i}\OperatorTok{);}
    \OperatorTok{\}}
\NormalTok{    printf}\OperatorTok{(}\StringTok{"I Love You More Than }\SpecialCharTok{\%d\textbackslash{}n}\StringTok{"}\OperatorTok{,}\NormalTok{ n}\OperatorTok{);}
\OperatorTok{\}}
\end{Highlighting}
\end{Shaded}

\slidefig{assets/1.2.2-p14.jpg}{0.8}{课件：小练习 1}

设最深层循环的执行次数（总共循环的次数）为 \(x\)。追踪 \texttt{i}
的变化：

{\rowcolors{2}{wdblue!5}{white}\arrayrulecolor{wdblue!35}
\begin{xltabular}{\linewidth}{>{\hsize=1.371\hsize\raggedright\arraybackslash}X>{\hsize=0.857\hsize\raggedright\arraybackslash}X>{\hsize=0.857\hsize\raggedright\arraybackslash}X>{\hsize=0.857\hsize\raggedright\arraybackslash}X>{\hsize=0.857\hsize\raggedright\arraybackslash}X>{\hsize=1.200\hsize\raggedright\arraybackslash}X}
\rowcolor{wdblue}\textcolor{white}{\bfseries 循环次数} & \textcolor{white}{\bfseries 第 1 次后} & \textcolor{white}{\bfseries 第 2 次后} & \textcolor{white}{\bfseries 第 3 次后} & \textcolor{white}{\bfseries \ldots\ldots{}} & \textcolor{white}{\bfseries 第 \(x\) 次后} \\
\endfirsthead
\texttt{i} 的值 & \(2\) & \(4\) & \(8\) & \ldots\ldots{} & \(2^x\) \\ \hline
\end{xltabular}\par\vspace{4pt}}

由循环条件可知，循环结束时刚好满足 \(2^x>n\)，解得
\(x=\log_2 n+1\)，所以

\[T(n)=O(x)=O(\log_2 n)\]

\begin{deepbox}

课件的 \(x=\log_2 n+1\) 是近似，精确值
\(x=\lfloor\log_2 n\rfloor+1\)（\(n=3000\) 时循环 12 次，最后打印
4096）。不影响数量级 \(O(\log_2 n)\)。

\end{deepbox}

\hypertarget{ux5c0fux7ec3ux4e60-2ux641cux7d22ux6570ux5b57ux578bux7231ux4f60}{%
\subsection{小练习
2：搜索数字型爱你}\label{ux5c0fux7ec3ux4e60-2ux641cux7d22ux6570ux5b57ux578bux7231ux4f60}}

\cueblock{\cuenote}{\relax{为什么同一个算法会有不同的时间复杂度？最好、最坏、平均怎么算？}}

\begin{Shaded}
\begin{Highlighting}[]
\CommentTok{//算法4— 搜索数字型爱你}
\DataTypeTok{void}\NormalTok{ loveYou}\OperatorTok{(}\DataTypeTok{int}\NormalTok{ flag}\OperatorTok{[],} \DataTypeTok{int}\NormalTok{ n}\OperatorTok{)} \OperatorTok{\{}    \CommentTok{//n 为问题规模}
\NormalTok{    printf}\OperatorTok{(}\StringTok{"I Am Iron Man}\SpecialCharTok{\textbackslash{}n}\StringTok{"}\OperatorTok{);}
    \ControlFlowTok{for}\OperatorTok{(}\DataTypeTok{int}\NormalTok{ i}\OperatorTok{=}\DecValTok{0}\OperatorTok{;}\NormalTok{ i}\OperatorTok{\textless{}}\NormalTok{n}\OperatorTok{;}\NormalTok{ i}\OperatorTok{++)\{}          \CommentTok{//从第一个元素开始查找}
        \ControlFlowTok{if}\OperatorTok{(}\NormalTok{flag}\OperatorTok{[}\NormalTok{i}\OperatorTok{]==}\NormalTok{n}\OperatorTok{)\{}              \CommentTok{//找到元素n}
\NormalTok{            printf}\OperatorTok{(}\StringTok{"I Love You }\SpecialCharTok{\%d\textbackslash{}n}\StringTok{"}\OperatorTok{,}\NormalTok{ n}\OperatorTok{);}
            \ControlFlowTok{break}\OperatorTok{;}                   \CommentTok{//找到后立即跳出循环}
        \OperatorTok{\}}
    \OperatorTok{\}}
\OperatorTok{\}}

\CommentTok{//flag 数组中乱序存放了 1\textasciitilde{}n 这些数}
\DataTypeTok{int}\NormalTok{ flag}\OperatorTok{[}\NormalTok{n}\OperatorTok{]=\{}\FloatTok{1.}\ErrorTok{..n}\OperatorTok{\};}
\NormalTok{loveYou}\OperatorTok{(}\NormalTok{flag}\OperatorTok{,}\NormalTok{ n}\OperatorTok{);}
\end{Highlighting}
\end{Shaded}

\slidefig{assets/1.2.2-p15.jpg}{0.8}{课件：小练习 2}

这个算法的循环次数\textbf{取决于元素 \(n\)
在数组的哪个位置}，也就是取决于输入数据：

{\rowcolors{2}{wdblue!5}{white}\arrayrulecolor{wdblue!35}
\begin{xltabular}{\linewidth}{>{\hsize=0.638\hsize\raggedright\arraybackslash}X>{\hsize=1.468\hsize\raggedright\arraybackslash}X>{\hsize=0.894\hsize\raggedright\arraybackslash}X}
\rowcolor{wdblue}\textcolor{white}{\bfseries 情况} & \textcolor{white}{\bfseries 分析} & \textcolor{white}{\bfseries 结果} \\
\endfirsthead
最好情况 & 元素 \(n\) 在第一个位置，循环 1 次 & \hlgreen{最好时间复杂度} \(T(n)=O(1)\) \\ \hline
最坏情况 & 元素 \(n\) 在最后一个位置，循环 \(n\) 次 & \hlred{最坏时间复杂度} \(T(n)=O(n)\) \\ \hline
平均情况 & 假设元素 \(n\) 在任意一个位置的概率相同，为 \(\frac{1}{n}\) & \hlblue{平均时间复杂度} \(T(n)=O(n)\) \\ \hline
\end{xltabular}\par\vspace{4pt}}

平均情况的计算（数学期望）：元素 \(n\) 在第 \(k\) 个位置时要循环 \(k\)
次，概率都是 \(\frac1n\)，所以

\[x = (1+2+3+\cdots+n)\cdot\frac{1}{n} = \frac{n(1+n)}{2}\cdot\frac{1}{n} = \frac{1+n}{2},\qquad T(n)=O(x)=O(n)\]

\hypertarget{ux4e09ux79cdux65f6ux95f4ux590dux6742ux5ea6}{%
\subsection{三种时间复杂度}\label{ux4e09ux79cdux65f6ux95f4ux590dux6742ux5ea6}}

\cueblock{\cuenote}{\relax{三种时间复杂度的定义？一般考虑哪种？}}

\begin{defbox}

\begin{itemize}
\tightlist
\item
  \textbf{最坏时间复杂度}：最坏情况下算法的时间复杂度；
\item
  \textbf{平均时间复杂度}：所有输入示例\hlred{等概率出现}的情况下，算法的期望运行时间；
\item
  \textbf{最好时间复杂度}：最好情况下算法的时间复杂度。
\end{itemize}

\end{defbox}

\begin{teacherbox}

很多算法的执行时间与输入的数据有关。一般只考虑\textbf{最坏}和\textbf{平均}时间复杂度。最好时间复杂度参考意义不大，因为评价一个算法时，主要是想看它在不好的情况下会不会出现运行时间过长的问题。

\end{teacherbox}

\begin{bookbox}

王道书的说法：通常以\textbf{最坏时间复杂度}作为算法效率的评价标准，因为它提供了运行时间的\textbf{确定性上界}。

书上也举了一个查找的例子，只是方向反过来（从后往前找）。为了和课件对上，这里换成课件的命名：在数组
\texttt{flag{[}0..n-1{]}} 中\textbf{从最后一个元素往前}找值为 \texttt{n}
的元素（书上原文用的是数组 \texttt{A} 和要找的值
\texttt{k}，意思完全一样）。若 \texttt{flag{[}n-1{]}} 就是
\texttt{n}，基本操作 \texttt{i-\/-}
一次都不用执行（最好）；若数组里根本没有 \texttt{n}，要执行 \(n\)
次（最坏）。

\end{bookbox}

\hypertarget{ux516dux738bux9053ux4e66ux5f52ux7eb3ux603bux7ed3ux4e24ux7c7bux5e38ux89c1ux9898ux578bux7684ux5206ux6790ux65b9ux6cd5}{%
\section{\texorpdfstring{六、王道书“归纳总结”：两类常见题型的分析方法}{六、王道书归纳总结：两类常见题型的分析方法}}\label{ux516dux738bux9053ux4e66ux5f52ux7eb3ux603bux7ed3ux4e24ux7c7bux5e38ux89c1ux9898ux578bux7684ux5206ux6790ux65b9ux6cd5}}

\tocback

王道书章末“归纳总结”把时间复杂度题归为两类，课上没细讲，补在这里。

\hypertarget{ux7c7bux578b-1ux5faaux73afux4e3bux4f53ux4e2dux7684ux53d8ux91cfux53c2ux4e0eux5faaux73afux6761ux4ef6ux7684ux5224ux65ad}{%
\subsection{类型
1：循环主体中的变量参与循环条件的判断}\label{ux7c7bux578b-1ux5faaux73afux4e3bux4f53ux4e2dux7684ux53d8ux91cfux53c2ux4e0eux5faaux73afux6761ux4ef6ux7684ux5224ux65ad}}

\Needspace{4\baselineskip}

\cuenote{循环变量参与判断时，怎么列式求解？}\textbf{方法}：设基本运算的执行次数为
\(t\)，建立 \(t\) 与 \(n\) 之间的关系式，解出 \(t=f(n)\)。若 \(f(n)\)
的最高次幂为 \(k\)，则时间复杂度为 \(O(n^k)\)（对数、根号同理）。

\begin{Shaded}
\begin{Highlighting}[]
\CommentTok{// 例 1}
\DataTypeTok{int}\NormalTok{ i}\OperatorTok{=}\DecValTok{1}\OperatorTok{;}
\ControlFlowTok{while}\OperatorTok{(}\NormalTok{i}\OperatorTok{\textless{}=}\NormalTok{n}\OperatorTok{)}
\NormalTok{    i}\OperatorTok{=}\NormalTok{i}\OperatorTok{*}\DecValTok{2}\OperatorTok{;}

\CommentTok{// 例 2}
\DataTypeTok{int}\NormalTok{ y}\OperatorTok{=}\DecValTok{5}\OperatorTok{;}
\ControlFlowTok{while}\OperatorTok{((}\NormalTok{y}\OperatorTok{+}\DecValTok{1}\OperatorTok{)*(}\NormalTok{y}\OperatorTok{+}\DecValTok{1}\OperatorTok{)\textless{}}\NormalTok{n}\OperatorTok{)}
\NormalTok{    y}\OperatorTok{=}\NormalTok{y}\OperatorTok{+}\DecValTok{1}\OperatorTok{;}
\end{Highlighting}
\end{Shaded}

\begin{itemize}
\tightlist
\item
  \textbf{例 1}：设 \texttt{i=i*2} 执行 \(t\) 次，则 \(2^t\le n\)，解得
  \(t\le\log_2 n\)，故 \(T(n)=O(\log_2 n)\)。（和小练习 1 一样）
\item
  \textbf{例 2}：设 \texttt{y=y+1} 执行 \(t\) 次，则 \(t=y-5\)，且
  \((t+5+1)(t+5+1)<n\)，解得 \(t<\sqrt{n}-6\)，故 \(T(n)=O(\sqrt{n})\)。
\end{itemize}

\hypertarget{ux7c7bux578b-2ux5faaux73afux4e3bux4f53ux4e2dux7684ux53d8ux91cfux4e0eux5faaux73afux6761ux4ef6ux65e0ux5173}{%
\subsection{类型
2：循环主体中的变量与循环条件无关}\label{ux7c7bux578b-2ux5faaux73afux4e3bux4f53ux4e2dux7684ux53d8ux91cfux4e0eux5faaux73afux6761ux4ef6ux65e0ux5173}}

\Needspace{4\baselineskip}

\cuenote{递归程序的时间复杂度怎么求？}此类问题可用\textbf{数学归纳法}或\textbf{直接累计循环执行次数}。多层循环从内向外逐层计算，忽略单步语句、条件判断语句，只关注主体语句的执行次数。

\begin{itemize}
\tightlist
\item
  \textbf{递归程序}：通常建立递推关系式求解。例如每次调用做常数次操作、再递归调用规模为
  \(n-1\) 的自己：
\end{itemize}

\[T(n)=1+T(n-1)=1+1+T(n-2)=\cdots=(n-1)+T(1),\qquad T(n)=O(n)\]

\begin{itemize}
\tightlist
\item
  \textbf{非递归程序}：直接累加次数即可。
\end{itemize}

\hypertarget{ux4e03ux5c0fux6545ux4e8bux4e3aux4ec0ux4e48ux53eaux5173ux5fc3-n-ux8db3ux591fux5927ux7684ux60c5ux51b5}{%
\section{\texorpdfstring{七、小故事：为什么只关心 \(n\)
足够大的情况}{七、小故事：为什么只关心 n 足够大的情况}}\label{ux4e03ux5c0fux6545ux4e8bux4e3aux4ec0ux4e48ux53eaux5173ux5fc3-n-ux8db3ux591fux5927ux7684ux60c5ux51b5}}

\tocback

\cueblock{\cuenote}{\relax{为什么时间复杂度只关注 \(n\) 很大的情况？}}

\begin{teacherbox}

很多算法的性能问题，只有在问题规模足够大的时候才会暴露出来。我工作时写过一个玩家匹配对战的算法，用了两层嵌套循环，时间复杂度
\(O(n^2)\)。自己测试只用了两三百个玩家，逻辑没问题。部署到线上后，同时在线用户有几十万，算法一跑，几十万玩家都受到了影响。后来师傅帮我优化成了
\(O(n)\)。

\end{teacherbox}

{\rowcolors{2}{wdblue!5}{white}\arrayrulecolor{wdblue!35}
\begin{xltabular}{\linewidth}{>{\hsize=0.957\hsize\raggedright\arraybackslash}X>{\hsize=1.021\hsize\raggedright\arraybackslash}X>{\hsize=1.021\hsize\raggedright\arraybackslash}X}
\rowcolor{wdblue}\textcolor{white}{\bfseries 规模} & \textcolor{white}{\bfseries \(O(n^2)\) 的基本操作次数} & \textcolor{white}{\bfseries \(O(n)\) 的基本操作次数} \\
\endfirsthead
本地测试 \(n\approx 100\) & \(10^4\)（一万） & \(10^2\) \\ \hline
线上 \(n=10^5\)（十万） & \(10^{10}\)（一百亿） & \(10^5\)（十万） \\ \hline
\end{xltabular}\par\vspace{4pt}}

\begin{fixbox}

字幕里老师说优化后“时间开销也就是 1 万这么多”，按 \(n=10^5\)、\(O(n)\)
算应为 \(10^5\)（十万），口误不影响结论：\(10^5\) 比 \(10^{10}\)
小了整整 \(10^5\) 倍。

\end{fixbox}

\hypertarget{ux516bux77e5ux8bc6ux56deux987eux62ffux5230ux4e00ux6bb5ux4ee3ux7801ux600eux4e48ux6c42-tn}{%
\section{\texorpdfstring{八、知识回顾：拿到一段代码怎么求
\(T(n)\)}{八、知识回顾：拿到一段代码怎么求 T(n)}}\label{ux516bux77e5ux8bc6ux56deux987eux62ffux5230ux4e00ux6bb5ux4ee3ux7801ux600eux4e48ux6c42-tn}}

\tocback

\slidefig{assets/1.2.2-p17.jpg}{0.8}{课件：本节知识回顾}

\Needspace{4\baselineskip}

\cuenote{求时间复杂度的三步是什么？两条规则？口诀？}这一节的内容是一套“方法”，所以回顾也按做题时的思考顺序排：先判断代码长什么样，再套对应的做法。课件上的结论都放在了对应的那一行里。

{\rowcolors{2}{wdblue!5}{white}\arrayrulecolor{wdblue!35}
\begin{xltabular}{\linewidth}{>{\hsize=0.780\hsize\raggedright\arraybackslash}X>{\hsize=1.440\hsize\raggedright\arraybackslash}X>{\hsize=0.780\hsize\raggedright\arraybackslash}X}
\rowcolor{wdblue}\textcolor{white}{\bfseries 看到的代码} & \textcolor{white}{\bfseries 怎么做} & \textcolor{white}{\bfseries 结果} \\
\endfirsthead
全是顺序语句，没有循环 & 顺序执行的代码只影响常数项（结论 1） & \(O(1)\) \\ \hline
一层循环，变量每次 +1 & 挑循环里的一条基本操作，数它的执行次数（结论 2） & \(O(n)\) \\ \hline
循环变量每次 ×2 & 设执行 \(x\) 次，由 \(2^x>n\) 解出 \(x\)（小练习 1） & \(O(\log_2 n)\) \\ \hline
条件里是 \(y^2<n\) 之类 & 设执行 \(t\) 次，代入条件解出 \(t\)（书上类型 1） & 如 \(O(\sqrt n)\) \\ \hline
多层嵌套循环 & 只看最深层，各层次数相乘（结论 3、乘法规则） & 如 \(O(n^2)\) \\ \hline
几段代码并列 & 各自求完再相加，按“常对幂指阶”只留最高阶（加法规则） & 最高阶那项 \\ \hline
递归 & 列递推式，如 \(T(n)=1+T(n-1)\)（书上类型 2） & 如 \(O(n)\) \\ \hline
执行次数依赖输入数据 & 分最好、最坏、平均；一般只看最坏和平均（小练习 2） & 分情况写 \\ \hline
\end{xltabular}\par\vspace{4pt}}

\begin{exambox}

\begin{enumerate}
\def\labelenumi{\arabic{enumi}.}
\tightlist
\item
  \textbf{如何计算}：① 找一个基本操作（最深层循环）；② 分析其执行次数
  \(x\) 与 \(n\) 的关系 \(x=f(n)\)；③ \(x\) 的数量级 \(O(x)\) 就是
  \(T(n)\)。
\item
  \textbf{加法规则}：\(O(f(n))+O(g(n))=O(\max(f(n),g(n)))\)；\textbf{乘法规则}：\(O(f(n))\times O(g(n))=O(f(n)\times g(n))\)。
\item
  \textbf{常对幂指阶}：\(O(1)<O(\log_2 n)<O(n)<O(n\log_2 n)<O(n^2)<O(n^3)<O(2^n)<O(n!)<O(n^n)\)。
\item
  \textbf{三种复杂度}：最坏（输入数据“最坏”）、平均（所有输入等概率出现）、最好（输入数据“最好”）；一般考虑最坏和平均。
\end{enumerate}

\end{exambox}

\begin{pitfallbox}

\begin{itemize}
\tightlist
\item
  时间复杂度 \(O(n^2)\) 表示“执行时间与 \(n^2\)
  \textbf{成正比}（同数量级）“，不是”执行时间等于
  \(n^2\)“，也不是”问题规模是 \(n^2\)”。
\item
  比较复杂度时看\textbf{最高阶}，不看系数：\(10000\log_2 n\)
  的复杂度仍然比 \(n\) 低，哪怕系数大到一万。
\item
  算法的实际运行时间还受\textbf{输入数据初始状态}影响，所以才有最好、最坏、平均之分。
\end{itemize}

\end{pitfallbox}

\begin{summarybox}

\textbf{时间复杂度 = 基本操作执行次数随问题规模 \(n\) 增长的数量级。}

事后计时受机器、语言、编译器影响，还可能来不及测，所以改为事前预估
\(T(n)\) 与 \(n\) 的关系。\(n\) 足够大时，低阶项和系数都不重要，于是用大
O
只保留最高阶项。代码再长也不用逐行数：顺序语句只影响常数，嵌套循环只看最深层，挑一个基本操作数出它的执行次数
\(x=f(n)\)，\(O(x)\)
就是答案。判断谁的阶高背“常对幂指阶”；循环变量翻倍是
\(O(\log_2 n)\)；执行时间依赖输入的算法要分最好、最坏、平均，平均情况用期望计算。

\textbf{自测}：① \texttt{for(i=1;i\textless{}=n;i++)} 中
\texttt{i\textless{}=n} 执行几次？②
\texttt{i=1;\ while(i\textless{}=n)\ i=i*3;} 的时间复杂度？③
为什么最深层循环的执行次数一定最多？

\end{summarybox}
\end{document}
